\documentclass[11pt]{article}
\usepackage{amsmath,amssymb,amsthm,mathtools,booktabs,enumitem,geometry}
\geometry{margin=1in}
\title{Step 115: Zeta-Factorization Attempt for Shifted Boundary Packets}
\author{RH Membrane Construction Notes}
\date{}

\newtheorem{theorem}{Theorem}
\newtheorem{lemma}{Lemma}
\newtheorem{proposition}{Proposition}
\newtheorem{definition}{Definition}
\newtheorem{warning}{Warning}
\newtheorem{gate}{Gate}
\newcommand{\B}{\mathcal B}
\newcommand{\E}{\mathcal E}
\newcommand{\M}{\mathcal M}
\newcommand{\K}{\mathcal K}
\newcommand{\Ran}{\operatorname{Ran}}
\newcommand{\Ker}{\operatorname{Ker}}
\newcommand{\proj}{\operatorname{proj}}

\begin{document}
\maketitle

\section{Purpose}
Step 114 reduced the Burnol/co-Poisson visibility question to the inclusion
\[
  \Ran \mathfrak B_{\ell,a}\subseteq P_a=Y_a^\perp,
\]
where
\[
  \mathfrak B_{\ell,a}=J_aP_\infty\tau_\ell(I-P_\infty)
\]
is the shifted Sonin/prolate boundary packet transported to the Burnol carrier.  This note tests the strongest possible easy certificate for that inclusion: a zeta-factorization of every shifted boundary packet.

The target certificate is
\[
  \M(\mathfrak B_{\ell,a}u)(s)=\zeta(s)\alpha_{\ell,u}(s),
\]
with the correct Burnol support, pole, endpoint, and trivial-zero records.  If this holds, all nontrivial zero-evaluators annihilate the packet, and the packet lies in the co-Poisson complement.

The outcome is not a proof of factorization.  The result is a precise reduction and a residual: the raw shift mode contributes only a nonvanishing exponential Mellin multiplier, so a \(\zeta\)-factor cannot be created by the shift itself.  Factorization must be supplied by the surrounding projection/synthesis structure, or the residual must be absorbed by Hecke/Dirichlet source coercivity.

\section{Burnol carrier and zero-evaluator map}
Let \(L_a\) be Burnol's extended Sonine space, \(Y_a\subset L_a\) the closed span of the nontrivial zeta-zero evaluator vectors, and
\[
  P_a=Y_a^\perp
\]
the Burnol co-Poisson subspace for \(a<1\).  Let
\[
  \E_a:L_a\longrightarrow \mathcal Z_a
\]
denote the zero-evaluator map
\[
  (\E_af)_{\rho,k}=\M(f)^{(k)}(\rho),
  \qquad 0\le k<m_\rho.
\]
Then
\[
  P_a=\Ker \E_a.
\]
For finite zero windows \(Z_T\), write \(\E_{a,T}\) for the truncated evaluator map.

\begin{definition}[Boundary--co-Poisson residual]
For a shifted boundary block \(\mathfrak B_{\ell,a}\), define
\[
  \Xi^{\rm BC}_{\ell,a}=\mathfrak B_{\ell,a}^*\E_a^*\E_a\mathfrak B_{\ell,a}\succeq0,
\]
and its finite-window version
\[
  \Xi^{\rm BC}_{\ell,a,T}=\mathfrak B_{\ell,a}^*\E_{a,T}^*\E_{a,T}\mathfrak B_{\ell,a}.
\]
\end{definition}

\section{Factorization criterion}
\begin{theorem}[Zeta-factorization implies co-Poisson inclusion]
Suppose that for every legal input \(u\) there exists a meromorphic/Hardy-admissible function \(\alpha_{\ell,u}\) such that
\[
  \M(\mathfrak B_{\ell,a}u)(s)=\zeta(s)\alpha_{\ell,u}(s),
\]
with enough regularity to evaluate at each nontrivial zero \(\rho\) to order \(m_\rho-1\).  Then
\[
  \Ran\mathfrak B_{\ell,a}\subseteq P_a.
\]
Equivalently,
\[
  \Xi^{\rm BC}_{\ell,a}=0.
\]
\end{theorem}

\begin{proof}
For every nontrivial zero \(\rho\) of multiplicity \(m_\rho\), the product \(\zeta(s)\alpha_{\ell,u}(s)\) vanishes at \(s=\rho\) to order at least \(m_\rho\), provided the declared pole records do not cancel the zero.  Thus
\[
  \M(\mathfrak B_{\ell,a}u)^{(k)}(\rho)=0,
  \qquad 0\le k<m_\rho.
\]
Hence \(\E_a\mathfrak B_{\ell,a}u=0\) for every legal \(u\), so \(\Ran\mathfrak B_{\ell,a}\subseteq\Ker\E_a=P_a\).  The residual identity is immediate.
\end{proof}

\begin{theorem}[Inclusion is equivalent to zero-evaluator annihilation]
For any bounded legal boundary block \(\mathfrak B_{\ell,a}\),
\[
  \Ran\mathfrak B_{\ell,a}\subseteq P_a
  \quad\Longleftrightarrow\quad
  \E_a\mathfrak B_{\ell,a}=0
  \quad\Longleftrightarrow\quad
  \Xi^{\rm BC}_{\ell,a}=0.
\]
\end{theorem}

\begin{proof}
Since \(P_a=\Ker\E_a\), the first two statements are identical.  The third is equivalent because
\[
  \langle u,\Xi^{\rm BC}_{\ell,a}u\rangle=\|\E_a\mathfrak B_{\ell,a}u\|^2.
\]
\end{proof}

\section{The log-shift does not create the zeta factor}
Let \(F(y)\) be the logarithmic model of a multiplicative function, with Mellin/Fourier convention
\[
  \M F(s)=\int_\mathbb R F(y)e^{(1/2-s)y}\,dy.
\]
For a shift \((\tau_\ell F)(y)=F(y-\ell)\),
\[
  \M(\tau_\ell F)(s)=e^{(1/2-s)\ell}\M F(s).
\]
The multiplier \(e^{(1/2-s)\ell}\) is entire and nonvanishing.  Therefore the shift itself cannot introduce zeros at the nontrivial zeros of \(\zeta\).

\begin{warning}[No automatic factorization from shifts]
The raw finite-place log-shift modes \(\tau_{\log p}\) cannot by themselves produce a factor \(\zeta(s)\).  Any factorization
\[
  \M(P_\infty\tau_\ell(I-P_\infty)u)(s)=\zeta(s)\alpha_{\ell,u}(s)
\]
must come from the surrounding Sonin/prolate projection geometry, co-Poisson synthesis, or an additional source/descent record.  It is not a consequence of multiplicative translation.
\end{warning}

\section{Attempt status for one prime shift}
For \(\ell=\log p\), the finite local-factor expansion contributes shifted modes
\[
  \tau_{k\log p},\qquad k\in\mathbb Z\setminus\{0\}.
\]
Testing the primitive mode \(\ell=\log p\), the strongest factorization attempt would require
\[
  \E_aJ_aP_\infty\tau_{\log p}(I-P_\infty)=0.
\]
The available structure gives:
\begin{itemize}[leftmargin=2em]
  \item Burnol co-Poisson atoms have Mellin transforms with a zeta factor, hence lie in \(P_a\).
  \item The raw log-shift has only a nonvanishing exponential Mellin multiplier.
  \item The archimedean Sonin/prolate projection is a support/Fourier-cutoff object, not visibly a co-Poisson synthesis map.
  \item Step 104 found noncompact boundary packets under nonzero shifts; Step 105 showed compact repair cannot remove their essential class.
\end{itemize}
Thus the factorization route is not closed by existing structure.  The remaining exact object is
\[
  \Xi^{\rm BC}_{\log p,a}=\mathfrak B_{\log p,a}^*\E_a^*\E_a\mathfrak B_{\log p,a}.
\]

\section{Residual source absorption}
If the factorization route fails, the source route must absorb the residual.  The required theorem shape is:
\[
  \Xi^{\rm BC}_{S,a,+}\preceq F_n+E_n,
\]
where
\[
  F_n=\sum_{\omega\in\mathcal X_n}\lambda_{\omega,n}Q_\omega^*\Theta_\omega^{-1}Q_\omega.
\]
If additionally
\[
  F_n\succeq \Lambda_n(\Theta_0^-)^{-1},
  \qquad \Lambda_n\to\infty,
\]
and the defects vanish in the fixed/exhaustive ledger sense, then the boundary residual is source-absorbed.

\section{All-six record for the factorization gate}
\begin{center}
\begin{tabular}{lll}
\toprule
Channel & Required record & Status after Step 115\\
\midrule
P1 rewrite & log-shift to Mellin multiplier & exact, nonzero multiplier only\\
P2 feasibility & Burnol support/pole/null legality & declared, not fully discharged\\
P3 route & Sonin shift route vs co-Poisson route & obstructed unless annihilation holds\\
P4 staging & finite zero windows with tail & needed for finite tests\\
P5 packaging & Burnol \(L_a,Y_a,P_a\) carrier & accepted as target carrier\\
P6 audit & \(\Xi^{\rm BC}\) residual/source absorption & active gate\\
\bottomrule
\end{tabular}
\end{center}

\section{Conclusion}
The zeta-factorization route for a raw shifted boundary packet is not earned.  It reduces to the exact annihilation condition
\[
  \E_a\mathfrak B_{\log p,a}=0.
\]
Since the shift supplies only a nonvanishing exponential Mellin multiplier, any \(\zeta\)-factor must come from co-Poisson synthesis or a deeper semilocal projection identity.  In its absence, the load-bearing target becomes source absorption of
\[
  \Xi^{\rm BC}_{\log p,a}.
\]

\end{document}
